\documentclass[12pt,a4paper]{article}
\usepackage[utf8]{inputenc}
\usepackage{amsmath, amssymb}
\usepackage{geometry}
\usepackage{setspace}
\usepackage{graphicx}
\usepackage{enumitem}

\geometry{margin=2.5cm}
\setstretch{1.2}

\title{Taller de factorización, homogenización y racionalización}
\author{José Luciano Mejía Arias}
\date{}

\makeatletter
\renewcommand{\maketitle}{
  \begin{flushright}
    \includegraphics[width=3cm]{../../logo_unal.png}
  \end{flushright}
  \begin{center}
    {\LARGE \textbf{\@title}}\\[0.3cm]
    {\large \@author}
  \end{center}
}
\makeatother

\newcommand{\pts}[1]{\hfill\textbf{[#1 pts]}}

\begin{document}
\maketitle

\begin{center}
  {\large Valor total: 50 puntos}
\end{center}

\noindent Nombre: \rule{8cm}{0.4pt} \hfill Fecha: \rule{3cm}{0.4pt}

\vspace{0.5em}
\noindent\textit{Presente el desarrollo completo de cada ejercicio y simplifique al máximo.}

\section*{Parte I. Factorización \hfill (18 pts)}
\begin{enumerate}[leftmargin=2em, itemsep=1.1em]
  \item Factoriza: $81a^2-100b^2$ \pts{2}
  \item Factoriza: $x^3+2x^2-9x-18$ \pts{3}
  \item Factoriza: $15x^3y-25x^2y^2+5xy$ \pts{2}
  \item Factoriza: $4x^2+20x+25-9y^2$ \pts{3}
  \item Factoriza: $6x^2-x-15$ \pts{3}
  \item Factoriza completamente: $x^4-81$ \pts{2}
  \item Factoriza: $27x^3+64y^3$ \pts{3}
\end{enumerate}

\section*{Parte II. Homogenización \hfill (18 pts)}
\begin{enumerate}[leftmargin=2em, itemsep=1.1em]
  \item Multiplica y simplifica: $\sqrt{3}\cdot\sqrt[3]{9}$ \pts{3}
  \item Lleva a índice común: $\sqrt{3}$ y $\sqrt[3]{5}$ \pts{2}
  \item Efectúa: $\dfrac{3}{x-2}-\dfrac{2}{x+2}$ \pts{3}
  \item Ordena de menor a mayor: $\sqrt[6]{120},\ \sqrt{5},\ \sqrt[3]{11}$ \pts{3}
  \item Efectúa y simplifica:
    $\dfrac{1}{x^2-1}+\dfrac{2}{x^2+2x+1}-\dfrac{1}{x-1}$ \pts{5}
  \item Divide: $\dfrac{\sqrt[4]{8}}{\sqrt{2}}$ \pts{2}
\end{enumerate}

\section*{Parte III. Racionalización \hfill (14 pts)}
\begin{enumerate}[leftmargin=2em, itemsep=1.1em]
  \item Racionaliza: $\dfrac{6}{\sqrt{5}-\sqrt{2}}$ \pts{2}
  \item Racionaliza: $\dfrac{10}{\sqrt{5}}$ \pts{1}
  \item Racionaliza: $\dfrac{1}{1+\sqrt{2}+\sqrt{3}}$ \pts{3}
  \item Racionaliza: $\dfrac{4}{\sqrt[3]{2}}$ \pts{2}
  \item Racionaliza: $\dfrac{2+\sqrt{3}}{2-\sqrt{3}}$ \pts{2}
  \item Racionaliza: $\dfrac{7}{3\sqrt{7}}$ \pts{2}
  \item Racionaliza: $\dfrac{1}{\sqrt{x+1}-\sqrt{x}}$ \pts{2}
\end{enumerate}

\end{document}
